\section{统一序列接口：tokenization、autoregression 与 training objective}

上一章说明了为什么需要多模态；本章回答“如何把不同信号送进 Transformer”。最容易犯的错误是把 token 等同于整数 ID。课堂使用的 token 更像一个\textbf{按时间或空间排列的 model input unit}：文本可以是离散 subword ID，图像可以是 patch embedding 或 VQ code，音频可以是 frame-level vector。统一的是序列接口，不一定是表示类型。

\subsection{混合序列的 globally autoregressive view}

假设一个样本被展开为混合序列
$$
\mathbf{x}=(x_1,x_2,\ldots,x_T),\qquad m_t\in\{\text{text},\text{vision},\text{audio}\},
$$
其中 $m_t$ 表示第 $t$ 个位置的 modality。全局自回归建模写成
$$
p(\mathbf{x})=\prod_{t=1}^{T}p\!\left(x_t\mid x_{<t},m_{\le t}\right).
$$
这里 $x_t$ 可以是离散 token，也可以是连续 latent；后一种情形中，右侧的“概率”需要由 Gaussian、diffusion 或其他连续目标实现，而不能机械套用 vocabulary softmax。

\lecturefigure{slide-010.jpg}{把 text、image patch 与 audio frame 放进同一因果序列}{官方 deck 第 10 页；课堂 00:06:31--00:07:26}

\begin{knowledgebox}{读图：统一的是顺序与条件关系}
图中的 T、V、A 分别代表 text token、visual patch 与 audio frame。先看底部顺序：不同 modality 可以交错出现，Transformer 在一个 causal context 中读取它们。再看顶部两条原则：tokenization across the board 负责把信号变成可排列单元，globally autoregressive generation 负责规定“当前位置依赖此前全部位置”。这并不要求三个 modality 使用同一种 decoder 或 loss。
\end{knowledgebox}

\begin{lstlisting}[caption={将多模态样本打包成一个带 modality 标签的序列}]
def pack_sequence(text_ids, image_patches, audio_frames):
    sequence = []
    sequence += [("text", token) for token in text_ids]
    sequence += [("vision", patch) for patch in image_patches]
    sequence += [("audio", frame) for frame in audio_frames]
    return add_positions_and_boundaries(sequence)
\end{lstlisting}

\begin{warningbox}{autoregressive 不等于所有输出都逐像素生成}
“Globally autoregressive”描述 modality block 之间及序列位置之间的因果顺序。一个 image block 内部仍可由 diffusion 一次并行去噪，audio 也可使用专门 codec。若把它误读为“每个像素都像文字一样从左到右预测”，就会错过 Transfusion 的核心设计。
\end{warningbox}

\subsection{Token 不一定离散：四类表示接口}

不同 modality 的统计结构不同，因此 tokenization 的目标也不同。文本 tokenizer 追求 compact vocabulary 与语义复用；图像 patch embedding 保留局部连续特征；VQ tokenizer 用 codebook 压缩成离散 ID；音频 frame/codec token 在时间分辨率、音质与序列长度之间折衷。

\lecturefigure{slide-011.jpg}{文本、图像、视频和音频的 tokenization 入口}{官方 deck 第 11 页；课堂 00:05:19--00:06:20}

\begin{knowledgebox}{读图：先比较“切什么”，再比较“丢什么”}
文本通常切 subword；图像切空间 patch；视频在 patch 之外还切时间；音频切 frame 或 codec unit。切分越细，序列越长、计算越贵；切分越粗，局部细节和时间变化越容易丢失。真正的设计变量不是“有没有 tokenizer”，而是单位粒度、是否离散、压缩率、可逆性以及它服务 understanding 还是 generation。
\end{knowledgebox}

\lecturefigure{slide-012.jpg}{连续 embedding、离散 code 与生成 latent 的表示谱系}{官方 deck 第 12 页；课堂 00:06:20--00:07:10}

\begin{knowledgebox}{术语消化：multimodal representation}
\begin{longtable}{L{0.20\textwidth}L{0.27\textwidth}L{0.42\textwidth}}
\toprule
术语 & 解决的问题 & 核心机制与课程关系 \\
\midrule
Patch embedding & 把高分辨率图像变成有限序列 & 将固定大小 patch 线性投影为连续向量，适合理解，但不直接给出像素级生成分布。 \\
VQ token & 让图像进入离散 LM vocabulary & 每个 patch 选择最近 codebook entry，接口统一但存在 quantization information loss。 \\
Continuous latent & 降低像素空间生成难度 & VAE/encoder 把图像压到连续空间，后续可用 diffusion 预测，通常更适合生成。 \\
Audio frame / codec token & 表示高频时间信号 & frame 保留连续声学特征；codec token 更离散、更紧凑，但音质与语义粒度受 codec 限制。 \\
Modality boundary & 防止模型混淆序列语义 & 特殊 token、位置编码或 type embedding 标记 modality 与 block 边界。 \\
\bottomrule
\end{longtable}
\end{knowledgebox}

\teachervoice{00:07:00--00:07:10，老师特别提醒，dense vector 也可以被称为 token。讲义采用“序列单元”这一宽定义，避免把后文 continuous image latent 错判为不属于 tokenization。}

\subsection{Multimodal input 不等于 multimodal output}

若模型只在 text position 上计算 loss，它可以学会根据图像回答问题，却没有被训练去还原 image/audio distribution。令 $\mathcal{T}$ 是 text-output positions，文本交叉熵为
$$
\mathcal{L}_{\text{text-only}}
=-\sum_{t\in\mathcal{T}}\log p_\theta(x_t\mid x_{<t}).
$$
图像和音频 token 在这里是 conditioning context，而不是被预测目标。

\lecturefigure{slide-013.jpg}{多模态输入、仅文本输出的训练目标}{官方 deck 第 13 页；课堂 00:07:33--00:07:55}

\begin{knowledgebox}{读图：看 loss 落在哪些位置}
底部序列同时有 T、V、A，但顶部气泡只指向 text output。模型可以在 self-attention 中使用视觉和音频信息，却只因文字预测误差而更新。这样容易获得 captioning、VQA、document understanding 等能力，但不会自动获得 image synthesis 或 speech generation。
\end{knowledgebox}

\begin{lstlisting}[caption={只对文本输出位置计算交叉熵}]
def text_only_loss(hidden, targets, modality):
    logits = text_head(hidden)
    mask = modality == "text"
    return cross_entropy(logits[mask], targets[mask])
\end{lstlisting}

\lecturefigure{slide-014.jpg}{课堂列举的 text-output multimodal model regime}{官方 deck 第 14 页；课堂 00:07:58--00:08:12}

\begin{warningbox}{模型名称不等于固定 capability contract}
slide 以 2026 年产品名称说明“多模态输入、文本输出”这一常见训练/交互模式，但产品会持续迭代，且公开接口不等于完整训练目标。本讲义只保留 regime 层面的结论：\textbf{conditioning modality 与 predicted modality 必须分别审计。}
\end{warningbox}

\subsection{Omni objective：每个 modality 都有自己的预测头}

若模型要生成 text、image 与 audio，就需要让非文本位置也产生可优化误差。一个抽象写法是
$$
\mathcal{L}_{\text{omni}}
=\lambda_T\mathcal{L}_{T}
+\lambda_V\mathcal{L}_{V}
+\lambda_A\mathcal{L}_{A},
$$
其中 $\mathcal{L}_{T}$ 常是 categorical cross-entropy，$\mathcal{L}_{V}$ 可以是 diffusion/flow matching，$\mathcal{L}_{A}$ 可以是 codec-token CE 或连续声学目标。$\lambda_m$ 不只是数学权重，还决定 optimization bandwidth：某个 modality 的 gradient 是否会淹没其他任务。

\lecturefigure{slide-015.jpg}{Omni model 对 text、vision 与 audio 输出共同训练}{官方 deck 第 15 页；课堂 00:08:12--00:09:16}

\begin{importantbox}{Omni model 的最小判据}
Omni model 至少应把多个 modality 作为可生成目标，而不只是输入条件。它可以共享一个 backbone，也可以拥有 modality-specific heads、losses 或参数；“omni”描述 capability coverage，不代表 architecture 必须完全 homogeneous。
\end{importantbox}

\begin{warningbox}{Joint training 不保证 positive transfer}
把 losses 相加只保证梯度同时存在，不保证任务互相帮助。不同 modality 可能共享高层语义，也可能在低层统计、序列长度和梯度尺度上冲突。后面的 modality gap、MoT 与 understanding-generation asymmetry 都是在处理这个未解决问题。
\end{warningbox}

\subsection{本章小结}

统一多模态序列需要同时回答 representation、ordering 与 objective。Token 是序列单元，不必离散；global autoregression 规定跨 block 的条件关系，不排斥 block 内 diffusion；multimodal input 与 multimodal output 由 loss mask 明确区分。到这里，我们有了统一接口，但还不知道它是否能像语言模型一样规模化。

